\subsection{Nonabelian promotion test for the congruence rows}

\begin{definition}[Promotion datum]
For a congruence product row $F_j$ a nonabelian denominator promotion is
the following datum:
$$
\mathscr W_j=(L_j,P_j,\mathcal M_j,W_j,\epsilon_j,\rho_j,N_j(a)).
$$
Here $L_j$ is the Lorentzian degree lattice,
$$
P_j\subset L_j,\qquad (\delta,\delta)=2\quad(\delta\in P_j),
$$
is the set of simple real roots,
$$
\mathcal M_j=\{x\in L_j\otimes\R\mid (x,\delta)\leq0
\hbox{ for every }\delta\in P_j\},
$$
$$
W_j=\langle s_\delta\mid \delta\in P_j\rangle,
\qquad
\epsilon_j:W_j\to\{\pm1\}
$$
is the reflection character,
$$
(\rho_j,\delta)=-1\qquad(\delta\in P_j),
$$
and $N_j(a)$ are the imaginary simple-root coefficients on the
Weyl-sum side.  A product formula supplies only the exponent function in
the Euler product.  It does not supply $\mathscr W_j$.
\end{definition}

\begin{definition}[Degree lattices and pairing]
For $t\in\Z_{>0}$ set
$$
M_t=U(4t)\oplus\langle2\rangle.
$$
Write an element of $M_t$ as $(n,l,m)$.  The pairing is
$$
((n,l,m),(n',l',m'))_t
=2ll'-4t(nm'+n'm).
$$
The exponential convention is
$$
e_t^{(n,l,m)}=q^nr^ls^{tm}.
$$
Thus
$$
((n,l,m),(n,l,m))_t=-2(4tnm-l^2).
$$
\end{definition}

\begin{lemma}[Arithmetic real-root check]
Set
$$
P^{(1)}=\{(0,-1,0),(1,1,0),(0,1,1)\}\subset M_1,
\qquad
\rho^{(1)}=(\tfrac12,\tfrac12,\tfrac12),
$$
and
$$
P^{(2)}=\{(0,-1,0),(1,1,0),(1,3,1),(0,1,1)\}\subset M_2,
\qquad
\rho^{(2)}=(\tfrac14,\tfrac12,\tfrac14).
$$
Then
$$
((\delta,\delta'))_{\delta,\delta'\in P^{(1)}})
=
\begin{pmatrix}
2&-2&-2\\
-2&2&-2\\
-2&-2&2
\end{pmatrix},
$$
$$
((\delta,\delta'))_{\delta,\delta'\in P^{(2)}})
=
\begin{pmatrix}
2&-2&-6&-2\\
-2&2&-2&-6\\
-6&-2&2&-2\\
-2&-6&-2&2
\end{pmatrix},
$$
and
$$
(\rho^{(1)},\delta)_1=-1\quad(\delta\in P^{(1)}),
\qquad
(\rho^{(2)},\delta)_2=-1\quad(\delta\in P^{(2)}).
$$
\end{lemma}

\begin{proof}
Substitution into
$$
((n,l,m),(n',l',m'))_t=2ll'-4t(nm'+n'm)
$$
gives the two displayed Gram matrices.  The same substitution gives
the Weyl-vector equations.  All displayed real roots have square $2$,
so these equations are exactly
$$
(\rho,\delta)=-(\delta,\delta)/2.
$$
\end{proof}

\begin{proposition}[Row-by-row promotion status]
For the four congruence rows of Clery--Gritsenko \cite{GC}, the
nonabelian promotion test gives:
$$
\begin{array}{c|c|c|c|c}
j&F_j&\mathrm{CG\ row}&\mathrm{CHL/WKB\ host}&\mathrm{status}\\ \hline
5&\nabla_3&(t,N)=(1,2)&\mathcal G_1(2)&\mathrm{promoted}\\
6&\nabla_2&(t,N)=(1,3)&\mathcal G_1(3)&\mathrm{promoted}\\
7&\nabla_{3/2}&(t,N)=(1,4)&\mathcal G_1(4)&
\mathrm{promoted\ on\ the\ order\ }4\ \mathrm{cover}\\
8&Q_1&(t,N)=(2,2)&\mathcal G_2(4)&\mathrm{promoted}
\end{array}
$$
The data are:
$$
L_5=L_6=L_7=M_1,\qquad
P_5=P_6=P_7=P^{(1)},\qquad
\rho_5=\rho_6=\rho_7=\rho^{(1)},
$$
and
$$
L_8=M_2,\qquad P_8=P^{(2)},\qquad \rho_8=\rho^{(2)}.
$$
The chambers and Weyl groups are
$$
\mathcal M_j=\{x\in L_j\otimes\R\mid(x,\delta)\leq0
\hbox{ for all }\delta\in P_j\},
\qquad
W_j=\langle s_\delta\mid\delta\in P_j\rangle.
$$
All simple real roots are even, and
$$
\epsilon_j(w)=\det(w)\qquad(j=5,6,7,8).
$$
\end{proposition}

\begin{proof}
Clery--Gritsenko construct the automorphic products
$$
F_5=\nabla_3,\qquad F_6=\nabla_2,\qquad
F_7=\nabla_{3/2},\qquad F_8=Q_1
$$
as diagonal-divisor forms on the indicated congruence groups
\cite{GC}.  This gives the product exponents and divisors, but not the
Weyl-sum side.

Govindarajan's periodic table of twisted CHL denominator algebras
\cite{Gov11BKM} identifies
$$
\nabla_3\leftrightarrow\mathcal G_1(2),\qquad
\nabla_2\leftrightarrow\mathcal G_1(3),\qquad
\nabla_{3/2}\leftrightarrow\mathcal G_1(4),
$$
and
$$
Q_1\leftrightarrow\mathcal G_2(4).
$$
For fixed CHL level $N$, the algebras $\mathcal G_N(M)$ have the same
real simple roots, Weyl vector, and Cartan matrix; the imaginary simple
roots depend on $M$.  Thus the three $N=1$ rows use $P^{(1)}$ and
$A^{(1)}$, while $Q_1$ uses $P^{(2)}$ and $A^{(2)}$.
Cheng--Dabholkar identify the same real-root systems and prove the
Weyl-denominator interpretation for the CHL levels $N=1,2,3$
\cite{CD09BKM}.  The arithmetic lemma translates their real-root
matrices into the $M_t$ notation used here.

For $F_7$ the multiplier has order $4$ in the Clery--Gritsenko
classification.  The promotion is therefore a statement on the
$\mu_4$-multiplier cover.  It is not a statement that an ordinary
scalar product on $\Gamma_0^{(2)}(4)$ determines a denominator algebra.
\end{proof}

\begin{theorem}[Imported Weyl--Kac--Borcherds identities]
For $j=5,6,7,8$ there are integral imaginary simple-root coefficients
$N_j(a)$, imported from the CHL/WKB denominator theorem, such that
$$
\begin{aligned}
F_j(Z)
&=\sum_{w\in W_j}\det(w)
\left(
e^{w(\rho_j)}
+\sum_{a\in L_j\cap\R_{>0}\mathcal M_j}
N_j(a)e^{w(\rho_j+a)}
\right)\\
&=e^{\rho_j}
\prod_{\lambda>0}
\left(1-e^\lambda\right)^{\mu_j(\lambda)}.
\end{aligned}
$$
For $j=5,6,7$ the exponential is $e=e_1$; for $j=8$ it is $e=e_2$.
The product exponent $\mu_j(\lambda)$ is the sum of the
Clery--Gritsenko cusp-labelled exponents landing at $\lambda$.
\end{theorem}

\begin{proof}
The product side is Clery--Gritsenko's Borcherds product, with the
cusp-labelled exponent function.  The Weyl-sum side is not reconstructed
from that product.  It is the additional CHL theorem:
Govindarajan proves that the square-root dd-forms in the displayed
table are the Weyl--Kac--Borcherds denominator formulae of the
rank-three Lorentzian BKM Lie superalgebras $\mathcal G_N(M)$
\cite{Gov11BKM}.  Cheng--Dabholkar supply the underlying real-root and
Weyl-chamber construction for the relevant CHL levels \cite{CD09BKM}.
\end{proof}

\begin{remark}[What remains outside the promotion]
This promotion does not assert a denominator algebra from product data
alone.  It also does not construct a microscopic BPS Hilbert-space
bracket, and it does not build a row map from the Vol.~III physical
symbols $\Phi^{\mathrm{V3}}_j$ to the Clery--Gritsenko rows $F_j$.  In particular,
the promoted Clery--Gritsenko row
$$
F_7=\nabla_{3/2}
$$
must not be identified with a separate physical boundary row
$\Phi^{\mathrm{V3}}_7$ unless that row map has first been constructed.
\end{remark}
